\section{Conclusion}\label{sec:conclusion}

The apparent emerging-market gap in patent survival is mostly a matter of when the renewal clock starts and of what
gets patented. Counted from grant and at equal observable quality, Chinese and Korean firms keep and value their
patents much as firms in the United States, Japan, Germany and the United Kingdom do. Credit shocks and fee cuts
barely move their renewal decisions, and they do not drop home patents of inventions they still protect abroad.
For firms, the binding constraint is the invention, not the inventor.

Owner capacity matters at the margin: small emerging-market firms' patents survive much longer after a sale,
especially to another company. Policies that deepen markets for technology, such as registries, licensing
intermediaries and secure transfer of title, may therefore add value. Short-run credit or fee relief is unlikely to
change which patents are kept. The clear weak spot is university patenting in China, where patents are dropped
sooner and are worth well under half as much as comparable foreign ones, consistent with patenting driven by
evaluation and subsidy targets rather than use.

Three limitations stand out. ``Emerging markets'' here means mainly China and Korea; Brazil is too thin, and other
large emerging economies do not report lapses in PATSTAT. Citations do not capture every dimension of quality.
And the structural values assume standard fees, while small domestic owners in China and Korea can claim fee
reductions; bounding that effect is left to a revision. % TODO: fee-discount bounds; additional offices.
